\section{Problem Formulation and Claim Boundaries}
\label{sec:problem}

\subsection{Pool-Based Learning}

Let $x\in\mathbb{R}^{d}$ be a static PE feature vector and
$y\in\{0,1\}$ its benign/malware label. At acquisition round $r$, the labeled
and unlabeled portions of the development pool are $\mathcal{L}_r$ and
$\mathcal{U}_r$. Model $f_r$ is fitted only on $\mathcal{L}_r$ and assigns
probability-like score $p_i=f_r(x_i)$. Acquisition function $Q_r(x)$ ranks
$x\in\mathcal{U}_r$; the oracle reveals labels for the highest-ranked batch,
which is moved to $\mathcal{L}_{r+1}$.

Within each outer fold and seed, all policies share the same initial labeled set,
remaining pool, validation data, test data, model hyperparameters, and cumulative
budget. Test outcomes do not influence queries, target definition, threshold
selection, or stopping. The oracle is simulated by revealing retained labels. A
training-label budget is thus a count of oracle queries, not measured analyst
time or total dataset annotation cost.

\subsection{Learning-Curve Estimands}

Every run produces ROC--AUC and average precision (AP) at budgets
$B\in\{200,300,\ldots,1200\}$. The normalized area under a metric learning curve
$m(B)$ is
\begin{equation}
\mathrm{AULC}(m)=\frac{1}{B_{\max}-B_{\min}}
\int_{B_{\min}}^{B_{\max}}m(B)\,dB,
\label{eq:aulc}
\end{equation}
computed by trapezoidal integration over measured budgets. AULC summarizes how
quickly useful ranking is acquired; the endpoint is not discriminative when all
policies eventually train on the same full pool.

For each outer fold, a full-pool reference model defines a validation target
\begin{equation}
t_{\mathrm{fold}}=0.5+0.95\left(
\mathrm{AUC}^{\mathrm{full}}_{\mathrm{val}}-0.5\right).
\label{eq:target}
\end{equation}
Labels-to-target is the first training budget whose validation AUC reaches this
level. The full-pool reference and validation labels are additional experimental
resources and must be reported separately from queried training labels.

Low-FPR summaries include normalized partial AUC up to FPR 5\% and TPR learning
curve area at FPR caps of 1\% and 5\%. These quantities describe the measured
outer folds. They do not estimate enterprise false-alarm rates near
$10^{-3}$--$10^{-4}$, for which substantially more benign test samples are
required.

\subsection{Explanation Object}

For a fitted tree model and sample $x_i$, let
$\boldsymbol{\phi}_i=(\phi_{i1},\ldots,\phi_{id})$ be the signed TreeSHAP
contribution vector and $\phi_{i0}$ the bias. Native TreeSHAP is checked on the
raw-score scale:
\begin{equation}
f_{\mathrm{raw}}(x_i)=\phi_{i0}+\sum_{j=1}^{d}\phi_{ij}.
\label{eq:additivity}
\end{equation}
Numerical reconstruction verifies the implementation but not semantic feature
identity, causality, or analyst usefulness.

The deletion diagnostic selects sample-specific mask $M_i^{(k)}$ and replaces
selected coordinates by a training-derived baseline $b_j$:
\begin{equation}
\widetilde{x}_{ij}=\begin{cases}
b_j,&j\in M_i^{(k)},\\
x_{ij},&\text{otherwise}.
\end{cases}
\label{eq:mask}
\end{equation}
For top-SHAP deletion, $M_i^{(k)}$ contains the $k$ largest
$|\phi_{ij}|$. The control contains $k$ uniformly selected coordinates. Baseline
$b_j$ is either the training mean or median. This intervention measures fitted
model reliance in feature space; it need not correspond to a valid PE binary.

\subsection{Security Scope}

The measured label is generic Windows PE malware. Original SHA-256 identities
support traceability, but the study does not establish ransomware labels,
family-held-out performance, or zero-day detection. The adversarial scope also
excludes a white-box attacker who optimizes pool membership, labels, or features
against the acquisition rule. These limits are part of the estimand rather than
footnotes to a broader deployment claim.
